Probing review found defects in every package after the builders' gates and the council's reviews had run. This section reports what broke; Appendix~\ref{app:reviewstats} holds the review-record statistics and Appendix~\ref{app:builders} the builder record, whose \texttt{Assisted-by} trailers name \builderAgents{} coding agents led by \builderTopAgent{} (\builderTopAgentPrs{} of \buildersPrs{} pull requests). Every count below is a sourced minimum drawn from a public changelog, issue, grading note, or pull-request thread; Table~\ref{tab:defects} names the source per class, Appendix~\ref{app:locators} lists the locator behind every package cell, and Figure~\ref{fig:defects} charts the counts.

\paragraph{Probing review.}
The probing sessions ran separate from the builders, instructed to distrust the builders' claims, write their own attack policies, and verify against running code. The author configured each of those sessions and handed it a protocol; the reviewers were agent sessions, so this paper calls the pass a probing review, and the protocol the first-round sessions followed is not released. The released protocol in Appendix~\ref{app:protocol} is the re-review one, which the re-check sessions used after fixes landed. Two things about these sessions are not verifiable on the public record: authorship, because the sessions ran in coding-agent interfaces and the model behind them is not recorded in the repositories, and coverage, because a defect class nobody thought to probe would not appear. What the record does show is the output: findings, fixes, and locators that anyone can check again.

\paragraph{The taxonomy.}
Table~\ref{tab:defects} sorts the findings into \defectClasses{} classes, and Appendix~\ref{app:defectclasses} walks each class with its examples and sourced minima. The widest are strategies that ignore the task and still earn credit (at least \defectBlindStrategyCredit{} packages), honest answers scored zero over formatting alone (at least \defectHonestAnswerRejected{} packages), and hidden-split handling that was not hidden (every wave-1 package changed it at 0.2.0). The rest cover unsolvable tasks, unpunished violations, the SDK task-list cache, breakdowns the engine dropped, forged report rows, an evaluator a crafted answer could escape, documentation ahead of the code, generator bugs, lexical detectors that miss unseen wording, and prompts looser than the grader. All \reviewPackages{} packages needed at least one fix round before or after release, with a median of \reviewRoundsMedian{} rounds and a maximum of \reviewRoundsMax{} (data in \texttt{paper/data/review-rounds.csv}). The first \reviewWaveOne{} packages each shipped a 0.2.0 whose changelog states that 0.1.0 scores are not comparable; the fix briefs required that sentence, and the scoring changes behind it are real and on the same record.

\paragraph{The re-measurement.}
The E2 runs then re-measured the logged exploits on the shipped tarballs. \expExplRunRows{} of the \expExplFrozen{} frozen exploit rows produced clauses on published artifacts (the other \expExplPolicyOnly{} need the reviewer's unreleased policy), \expExplClauses{} clauses in all across single answers, episodes and source checkouts. Of the \expExplRecordClauses{} clauses that check a logged number, \expExplReplicated{} replicate on the shipped artifacts and \expExplNotReplicated{} do not; Table~\ref{tab:exploit-verdicts} lists each miss with its package, locator and measured value. On the attack rows the median score falls from \expExplAtkMedianOld{} at 0.1.0 to \expExplAtkMedianNew{} at 0.2.0, and the largest post-fix attack mean is \expExplAtkMaxNew{}. Of the \expExplPostfixClauses{} post-fix clauses, \expExplFixHolds{} hold and \expExplFixFails{} fail: the shipped 0.2.0 still gives the exploit credit on \expExplFixFailsPkgs{}, reported here as found rather than explained away. Figure~\ref{fig:exploits} plots the record against the re-measurement for every frozen row that produced a released-version number; the wave-2 rows below the rule have only a 0.1.0 release, some already carrying review fixes, so they are exploratory and sit outside the before/after claim. Because the exploit list is post hoc, the replication says the logged defects were real and mostly fixed on the shipped versions; it says nothing about how often probing misses a defect.

\paragraph{The review record, for and against.}
The record documents \approvalDefectCases{} approvals standing behind defects the probing review found later: pelican's pull request 4, issue 8, and sycophancy's pull request 20, approved while a publishing leak and a grading-note claim the code did not back were still in the tree. The counter-evidence is real too: the council approved \councilMergedApproved{} of the \councilMergedFeature{} merged feature pull requests and its reviews raised an open Critical or Major item on about \councilSeverityPrsRead{} of them, and one defect in Table~\ref{tab:defects}, the published hidden seed on sycophancy's pull request 20, was flagged by the council itself before merge.

\paragraph{What the record does and does not show.}
One caveat sits over the whole section: the defect table is built from the probing record, so ``probing found them'' is partly true by construction. The parts that do not rest on the record are the re-measurements: the E2 clause verdicts above, the E1 battery in Section~\ref{sec:prereg}, and the E4 contract probes in Section~\ref{sec:soundness}. Not every finding closed: crm-truthfulness ships as alpha with its lexical miss rate documented as a known limit, and pelican's structural check accepts drawings a person would reject and rejects some it would accept, a documented limit that is still open. A published mean that does not name its package version is also not reproducible, and Table~\ref{tab:suite} carries the versions for that reason.
